\begin{abstract}
Heterogeneous graph neural networks have achieved promising performance in modeling complex relational data with multiple types of nodes and edges.
However, they remain vulnerable to structural adversarial attacks, where a small number of malicious edge perturbations may be amplified through meta-path-based message propagation and further corrupt high-order semantic neighborhoods.
Existing robust heterogeneous graph methods mainly defend against such perturbations within the topological semantic space, for example by estimating edge confidence, filtering suspicious meta-path-induced edges, or learning semantic attention weights.
Although effective to some extent, these methods still heavily rely on the attacked heterogeneous topology to construct candidate neighborhoods, making it difficult to recover reliable representations when the topology itself has been severely polluted.

To address this limitation, we propose DVCL, a dual-view contrastive learning framework for robust heterogeneous graph learning.
DVCL constructs two complementary views: a purified topology view derived from meta-path-induced semantic structures, and a feature-induced view built from node feature similarity.
The topology view preserves heterogeneous structural semantics through meta-path modeling, semantic attention, and two-stage edge filtering, while the feature view provides a relatively topology-independent reference under structural attacks.
To further reduce semantic inconsistency between the two views, DVCL employs a symmetric cross-view contrastive objective that aligns the topology representation and feature representation of the same node while distinguishing representations from different nodes.
Together with supervised node classification and an auxiliary semantic topology learning objective, DVCL learns robust node representations that integrate heterogeneous semantic expressiveness and feature-level stability.
Experiments on benchmark heterogeneous graphs under structural attacks demonstrate that DVCL improves node classification robustness compared with existing heterogeneous graph defense baselines.
\end{abstract}

